\section{Processing Objective \& Analytical Question}

\subsubsection{Processing Objective}

The objective of this analysis is to use distributed PySpark processing to identify
transaction patterns and account relationships associated with money-laundering
activity within the anti-money laundering (AML) transaction dataset. The dataset contains a large volume of financial transactions, making it suitable for demonstrating distributed processing and
MapReduce-style analytical operations.

The analysis focuses on determining whether laundering activity is concentrated within
specific transaction characteristics, financial institutions, or repeated account
relationships. Specifically, the analysis examines transaction frequency, transaction
amounts, payment formats, transfer characteristics, originating banks, and sender-
receiver relationships. These analyses extend the exploratory findings from Part 1 by
applying distributed transformations, grouping operations and aggregations to extract
meaningful patterns from the large-scale transaction dataset.

\subsubsection{Central Question}

How can distributed PySpark processing be used to identify transaction patterns and
account relationships associated with money-laundering activity in the AML dataset?

To address this question, the analysis aims to:
\begin{enumerate}
    \item Identify transaction characteristics associated with higher observed laundering rates.
    \item Quantify laundering activity across originating banks and account relationships.
    \item Compare laundering behaviour across payment formats, transfer types and account relationships.
    \item Demonstrate how PySpark's distributed processing capabilities can perform large-scale analytical tasks through transformations, shuffle/grouping and reduction/aggregation operations.

\end{enumerate}

The resulting insights are intended to demonstrate how large-scale transaction analytics
can support the identification and prioritisation of potentially higher-risk transaction
patterns for further investigation.